\documentclass{article}

\usepackage{amssymb,amsmath,amsthm,mathtools,cancel}
\usepackage[margin=1in]{geometry}
\usepackage{enumitem}


\newtheorem{theorem}{Theorem}[section]
\newtheorem{lemma}[theorem]{Lemma}

\begin{document}
\noindent
\begin{minipage}[t]{0.68\textwidth}
{\large 6.1200 \textit{Mathematics for Computer Science}}\\[4pt]
{\large Problem Set 2}\\[4pt]
{\large Solutions by Harsh Bhalala}
\end{minipage}%
\hfill
\begin{minipage}[t]{0.27\textwidth}
\raggedleft
{\large 12, September 2026}
\end{minipage}

\vspace{6pt}

\noindent\rule{\textwidth}{0.5pt}

\section*{Problem 1. Fibonacci Divisibility}

\textbf{(a)}

\begin{proof}[Proof by strong induction]
Let $P(n) := F_{n + a} \equiv F_{n} \cdot F_{a+1} \pmod{F_a}$, now to prove $p(n)$ for $n \geq 0$ by strong induction.

\textbf{Base Case:}\\
\indent\indent $P(0) = F_{a} \equiv F_0 \cdot F_{a+1} \equiv 0 \pmod{F_a} $.\\
\indent\indent $P(1) = F_{a+1} \equiv F_1 \cdot F_{a+1} \equiv 1 \cdot F_{a+1} \equiv F_{a+1} \pmod{F_a}$\\
Therefore, both cases hold.

\textbf{Inductive step:} For strong induction on $n$ assume $P(n-1)$ and $P(n-2)$ holds for $n \geq 2$. So,
\begin{equation}
F_{(n-1)+a} \equiv F_{n-1} \cdot F_{a+1} \pmod{F_a}
\end{equation}
\begin{equation}
F_{(n-2)+a} \equiv F_{n-2} \cdot F_{a+1} \pmod{F_a}
\end{equation}
Now, by the properties of the congruence on equation (1) and (2),
\[
\begin{aligned}
F_{(n-1)+a} + F_{(n-2)+a} &\equiv F_{n-1} \cdot F_{a+1} + F_{n-2} \cdot F_{a+1} &\pmod{F_a}\\
F_{n+a} &\equiv \left(F_{n-1} + F_{n-2} \right) \cdot F_{a+1} &\pmod{F_a}\\
F_{n+a} &\equiv F_n \cdot F_{a+1} &\pmod{F_a}
\end{aligned}
\]
Which implies $P(n)$ is true. Therefore, by the strong induction hypothesis, $P(n)$ is true for all $n \geq 0$.
\end{proof}


\vspace{\baselineskip}
\noindent\textbf{(b)}

\begin{proof}[Proof by induction]
Induction hypothesis Assume \(P(n)\) is true, i.e.,

\[F_{na} \equiv 0 \pmod{F_a}.\]
\textbf{Base Case:} $P(0)= F_0 \equiv 0 \pmod{F_a}$

Base case holds. We now prove $P(n+1)$, namely,
\[F_{(n+1)a} \equiv 0 \pmod{F_a}.\]
\[F_{na+a} \equiv 0 \pmod{F_a}.\]

Using part (a) with $n$ replaced by $na$,
\[
\begin{aligned}
F_{na+a} &\equiv F_{na}F_{a+1} \pmod{F_a}. \\[6pt]
F_{na+a} &\equiv 0\cdot F_{a+1} \pmod{F_a}. \\[6pt]
F_{na+a} &\equiv 0 \pmod{F_a}.
\end{aligned}
\]
This proves $P(n+1)$ holds, and by the induction hypothesis $P(n)$ is true for $n \geq 0$.
\end{proof}
Now if $a \vert b$, then $\exists k \in \mathbb{N}$ such that $a * k = b$. Now by $P(k)$ , $F_{ak} = F_b \equiv 0 \pmod{F_a}$ which implies $F_a \vert F_b$.

\section*{Problem 2. Computing Modular Inverses}

\textbf{(a)}
First, using Euclid's algorithm:
\[
\begin{aligned}
gcd(43,15) &= gcd(15,13) \qquad 13 = 42 - 2 \cdot 15\\
&= gcd(13,2) \qquad 2 = 15 - 1 \cdot 13\\
&= gcd(2,1) \qquad 1 = 13 - 6 \cdot 2\\
\end{aligned}
\]
Now doing backwards substitution (Pulverizer):\\
\textbf{Substitute 2 (from $2 = 15 - 1 \cdot 13$)}
\[
\begin{aligned}
1 = 13 - 6 \cdot (15 - 1 \cdot 13)\\
1 = 7 \cdot 13 - 6 \cdot 15\\
\end{aligned}
\]
\textbf{Substitute 13 (from $13 = 43 - 2 \cdot 15$)}
\[
\begin{aligned}
1 = 7 \cdot (43 - 2 \cdot 15) - 6 \cdot 15\\
1 = 7 \cdot 43 - 14 \cdot 15 - 6 \cdot 15\\
1 = 7 \cdot 43 - 20 \cdot 15\\
\end{aligned}
\]
The coefficient of 15 is -20. Therefore, the inverse would be $43 - 20 = 23$. So $23 \cdot 15 \equiv 1 \pmod{43}$.

\vspace{\baselineskip}
\noindent\textbf{(b)}
From Fermat's Little theorem we know that $15^{42} \equiv 1 \pmod{43}$ and now multiplying both side with $15^{-1}$
we get
\[ 15^{41} \equiv 15^{-1} \pmod{43}\]
Now to simplify $15^{41}$ we caclculate the smaller powers of 15:
\[
\begin{aligned}
15^1 \equiv 15 \pmod{43}\\
15^2 = 225 \equiv 10 \pmod{43}\\
15^4 \equiv 100 \equiv 14 \pmod{43}\\
15^8 \equiv 196 \equiv 24 \pmod{43}\\
15^{16} \equiv 576 \equiv 17 \pmod{43}\\
15^{32} \equiv 289 \equiv 31 \pmod{43}\\
\end{aligned}
\]
Now $15^{41} = 15^{32} \cdot 15^{8} \cdot 15$ Therefore, $15^{-1} \equiv 31 \cdot 24 \cdot 15 \equiv 23 \pmod{43}$
Meaning the inverse of $15 \text{ mod } 43$ is 23.


\section*{Problem 3. Primality Testing}

\textbf{(a)}
\begin{proof}[Proof by Contradiction]
Assume an integer $1 \leq a \leq n - 1$ is fermat's witness of some prime number $n$.
since $a$ is fermat's witness $a^{n-1} \not\equiv 1 \pmod{n}$.
Also since $n$ is prime and $a$ is in range $[1,n-1]$, Therefore $a$ is not multiple of $n$ so
by fermat's little theorem $a^{n-1} \equiv 1 \pmod{n}$ which is a contradiction.
\end{proof}

\vspace{\baselineskip}
\noindent\textbf{(b)}
\begin{proof}
Assume $p$ is a prime and $x^2 \equiv_p 1$ by the defintion of congruence $p \mid x^2 - 1$
\[p \mid (x - 1)(x + 1)\]
Since $p$ is prime (by Lemma 9.4.2) $p \mid (x - 1)$ or $p \mid (x + 1)$ so by defintion of congruence $x \equiv_p -1$ or $x \equiv_p 1$.
\end{proof}
\pagebreak

\noindent\textbf{(c)}
\begin{proof}[Proof by contradiction]
Assume, for the sake of contradiction, that the algorithm returns 'composite' for some prime number $n$.
,So at least one of the conditions from the algorithm must hold.
\begin{itemize}
\item $x_e \not\equiv_n 1$
\item Consecutive terms $(x_i, x_{i+1})$ where $x_{i+1} \equiv_n 1$ but $x_i \not\equiv_n \pm 1$.
\end{itemize}

Suppose the first condition holds; then $x_e = a^{n-1} \not\equiv_n 1$.
Now Since $n$ is prime and $1 \leq a \leq n-1$ and
Fermat's Little Theorem implies that $a^{n-1} \equiv_n 1$, which is a contradiction.

The first condition led to a contradiction, so the second condition must be true.
So $x_{i+1} = {x_i}^2 \equiv_n 1$ but $x_i \not\equiv_n \pm 1$.
Since n is prime, by lemme from part (b), the given condition yields a contradiction.

Both conditions lead to a contradiction, meaning our assumption must be wrong, so n must indeed be composite.
\end{proof}

\section*{Problem 4. Pulverizer State Machine}
\textbf{(a)}
For some arbitrary state $(x, y, s, t, u, v)$ say Pres1, Pres2, Pres3 are true and we need to prove that Pres1, Pres2, Pres3 are true for the next state $(x', y', s', t', u', v')$.

By the definition of transition in the given state machine, we have $x' = y$, $y' = (x \text{ rem } y)$, $s' = u$, $t' = v$, $u' = s - qu$, and $v' = t - qv$, where $x = qy + (x \text{ rem } y)$.
\begin{itemize}

\item \textbf{Pres1:} gcd($x', y'$) = gcd($a, b$) we know that gcd($x', y'$) = gcd($y, x \text{ rem } y$) = gcd($x, y$) = gcd($a, b$) which proves Pres1 is preserved.

\item \textbf{Pres2:} $s'a + t'b = x'$, using the definition of $s', t', x'$ we have $s'a + t'b = ua + vb = y = x'$ which proves Pres2 is preserved.

\item \textbf{Pres3:} $u'a + v'b = y'$, using the definition of $u', v', y'$ we have
\[
\begin{aligned}
u'a + v'b &= (s - qu)a + (t - qv)b\\
&= (sa + tb) - q(ua + vb)\\
&= x - q(y) \\
&= (x \text{ rem } y)\\
&= y'\\
\end{aligned}
\]
which proves Pres3 is preserved.
\end{itemize}

\vspace{\baselineskip}
\noindent\textbf{(b)} From part (a), we proved that Pres1 and Pres2 are preserved, and from the definition of the given state machine we know that
In the final state, $y = 0$, so now substituting $y$ and Pres2 in Pres1, we get gcd($a, b$) = gcd($x, y$) = gcd(sa + tb, 0) = sa + tb.
Therefore, $s$ and $t$ are the coefficients satisfying B\'ezout's identity.

\vspace{\baselineskip}
\noindent\textbf{(c)} Because the given state machine is just the Euclidean algorithm with some extra computation, and even the termination condition is
same as the Euclidean algorithm, which is when y = 0.


\end{document}